\begin{apartats}

\apartat[1,25]{0,75}
Comprova que $F(x)=(x^2+1)^3$ és una primitiva de $f(x)=6x(x^2+1)^2$. Explica per què
$G(x)=(x^2+1)^3-5$ també n'és una primitiva.

\begin{solucio}
$F'(x)=3(x^2+1)^2\cdot2x=6x(x^2+1)^2=f(x)$ per a tot $x$, i per tant $F$ és una primitiva de $f$.\\
$G(x)=F(x)-5$, i la derivada d'una constant és 0: $G'(x)=F'(x)=f(x)$. Totes les funcions $F(x)+k$ són
primitives de $f$.
\end{solucio}

\begin{tria}{primitiva-condicions}
\itemtria{un-punt}{1}{1,25}
Troba la primitiva $F$ de la funció $f(x)=3x^2-4x+1$ que compleix $F(1)=2$.

\begin{solucio}
Les primitives de $f$ són $F(x)=x^3-2x^2+x+k$. Com que $F(1)=1-2+1+k=k$, la condició $F(1)=2$ dona
$k=2$: $F(x)=x^3-2x^2+x+2$.
\end{solucio}

\itemtria{segona-derivada}{1}{1,25}
Troba la funció $F$ que compleix $F''(x)=6x$, $F'(1)=2$ i $F(0)=1$.

\begin{solucio}
$F'$ és una primitiva de $F''$: $F'(x)=3x^2+k_1$, i $F'(1)=3+k_1=2$ dona $k_1=-1$.\\
$F$ és una primitiva de $F'(x)=3x^2-1$: $F(x)=x^3-x+k_2$, i $F(0)=k_2=1$. Per tant, $F(x)=x^3-x+1$.
\end{solucio}
\end{tria}

\begin{nomesllarg}
\apartat{0,75}
Les funcions $F(x)=\sin^2x$ i $G(x)=-\cos^2x$, són primitives de la mateixa funció? En quant es
diferencien?

\begin{solucio}
$F'(x)=2\sin x\cos x=\sin(2x)$ i $G'(x)=-2\cos x\cdot(-\sin x)=\sin(2x)$: totes dues són primitives de
$\sin(2x)$. Es diferencien en una constant: $F(x)-G(x)=\sin^2x+\cos^2x=1$.
\end{solucio}
\end{nomesllarg}

\end{apartats}
